% Part: lambda-calculus
% Chapter: church-rosser
% Section: parallel-beta-reduction

\documentclass[../../../include/open-logic-section]{subfiles}

\begin{document}

\olfileid{lam}{cr}{pb}

\olsection{Reduksi-$\beta$ paralel}

Kita ngenalake konsep
\emph{reduksi-$\beta$ paralel}
lan mbuktekake yen relasi iki
nduweni sipat Church--Rosser.

\begin{defn}[reduksi-$\beta$ paralel, $\bredpar$] \ollabel{defn:bredpar}
  Reduksi paralel ($\bredpar$)
  ing suku ditetepake kanthi
  induktif kaya mangkene:
  \begin{enumerate}
    \item \ollabel{defn:bredpar1} $x \bredpar x$.
    \item \ollabel{defn:bredpar2} Yen $N \bredpar N'$, mula $\lambd[x][N] \bredpar
      \lambd[x][N']$.
    \item \ollabel{defn:bredpar3} Yen $P \bredpar P'$ lan $Q \bredpar Q'$, mula $PQ \bredpar
      P'Q'$.
    \item \ollabel{defn:bredpar4} Yen $N \bredpar N'$ lan $Q \bredpar Q'$, mula
      $(\lambd[x][N])Q \bredpar \Subst{N'}{Q'}{x}$.
  \end{enumerate}
\end{defn}

Reduksi-$\beta$ paralel ngidini
pirang-pirang redex ing sawijining
suku dikontraksi sajroning siji
langkah. Bedane karo
reduksi-$\beta$, ing sawijining
langkah paralel kita mung bisa
ngontraksi redex kang wis ana
ing suku asal, dudu redex anyar
kang muncul saka reduksi-$\beta$
paralel.
Contone, suku
$(\lambd[f][fx])(\lambd[y][y])$
sajroning siji langkah reduksi-$\beta$
paralel
mung bisa direduksi dadi suku
iku dhewe utawa
$(\lambd[y][y])x$, durung
bisa dadi~$x$, sanajan suku
asal bisa direduksi-$\beta$
dadi~$x$. Redex kanggo langkah
pungkasan iku lagi muncul
sawisé siji langkah reduksi-$\beta$
paralel. Nanging langkah reduksi-$\beta$
paralel
kapindho ngasilake~$x$.

\begin{thm}\ollabel{thm:refl}
  $M \bredpar M$.
\end{thm}

\begin{proof}
  Latihan.
\end{proof}

\begin{prob}
  Buktekna \olref[lam][cr][pb]{thm:refl}.
\end{prob}

\begin{defn}[pangembangan lengkap-$\beta$]\ollabel{defn:bcd}
  \emph{Pangembangan lengkap-$\beta$}~$\bcd{M}$
  saka $M$ ditetepake kanthi
  induktif kaya mangkene:
  \begin{align}
    \bcd{x} &= x \ollabel{defn:bcd1} \\
    \bcd{(\lambd[x][N])} &= \lambd[x][\bcd{N}] \ollabel{defn:bcd2}\\
    \bcd{(PQ)} &= \bcd{P}\bcd{Q} && \text{yen $P$ dudu abstraksi-$\lambd$}
    \ollabel{defn:bcd3} \\
    \bcd{((\lambd[x][N])Q)} &= \Subst{\bcd{N}}{\bcd{Q}}{x} \ollabel{defn:bcd4}
  \end{align}
\end{defn}

Pangembangan lengkap-$\beta$
sawijining suku, kaya
kang katuduhake dening jenenge,
yaiku ``reduksi paralel kang
lengkap''. Ing reduksi-$\beta$
paralel kita isih bisa milih
ora ngontraksi sawijining redex,
nanging ing pangembangan lengkap
kabeh redex asal kudu dikontraksi.
Mula pangembangan lengkap saka
$(\lambd[f][fx])(\lambd[y][y])$
yaiku $(\lambd[y][y])x$,
dudu suku asal iku dhewe.

\begin{editorial}
  Definisi iki isih nduweni
  masalah: kita durung nerangake
  carane netepake fungsi ing
  suku-$\lambd$ kanthi rekursif.
  Bab iki bakal didandani
  ing tembe.
\end{editorial}

\begin{lem}\ollabel{lem:comp}
  Yen $M \bredpar M'$ lan $R \bredpar R'$,
  mula $\Subst{M}{R}{y}
  \bredpar \Subst{M'}{R'}{y}$.
\end{lem}

\begin{proof}
  Gunakake induksi ing !!{derivation}
  saka $M \bredpar M'$.
  \begin{enumerate}
    \item Yen langkah pungkasan
      \olref{defn:bredpar1}: latihan.
    \item Yen langkah pungkasan
      \olref{defn:bredpar2}: $M$
      awujud $\lambd[x][N]$ lan
      $M'$ awujud $\lambd[x][N']$,
      kanthi $N \bredpar N'$.
      Kita kudu mbuktekake
      $\Subst{(\lambd[x][N])}{R}{y} \bredpar
      \Subst{(\lambd[x][N'])}{R'}{y}$,
      yaiku
      $\lambd[x][\Subst{N}{R}{y}] \bredpar
      \lambd[x][\Subst{N'}{R'}{y}]$.
      Iki ndherek saka
      \olref{defn:bredpar2} lan
      hipotesis induksi, sawisé
      wakil binder dipilih seger.
    \item Yen langkah pungkasan
      \olref{defn:bredpar3}: latihan.
    \item Yen langkah pungkasan
      \olref{defn:bredpar4}: $M$
      awujud $(\lambd[x][N])Q$
      lan $M'$ awujud
      $\Subst{N'}{Q'}{x}$.
      Kita kudu mbuktekake
      $\Subst{((\lambd[x][N])Q)}{R}{y} \bredpar
      \Subst{\Subst{N'}{Q'}{x}}{R'}{y}$,
      yaiku
      $(\lambd[x][\Subst{N}{R}{y}])\Subst{Q}{R}{y} \bredpar
      \Subst{\Subst{N'}{R'}{y}}{\Subst{Q'}{R'}{y}}{x}$.
      Langkah iki nggunakake
      \olref{defn:bredpar4},
      hipotesis induksi, lan
      hukum komposisi substitusi
      kanthi binder kang dipilih
      seger. Justifikasi pungkasan
      iki durung diwenehake ing
      teks sumber.
  \end{enumerate}
\end{proof}

\begin{prob}
  Rampungna bukti \olref[lam][cr][pb]{lem:comp}.
\end{prob}

\begin{lem}\ollabel{lem:cont}
  Yen $M \bredpar M'$, mula
  $M' \bredpar \bcd{M}$.
\end{lem}
\begin{proof}
  Gunakake induksi ing !!{derivation}
  saka $M \bredpar M'$.
  \begin{enumerate}
    \item Yen aturan pungkasan
      \olref{defn:bredpar1}: latihan.
    \item Yen aturan pungkasan
      \olref{defn:bredpar2}: $M$
      awujud $\lambd[x][N]$
      lan $M'$ awujud
      $\lambd[x][N']$ kanthi
      $N \bredpar N'$.
      Kita kudu nuduhake
      $\lambd[x][N'] \bredpar
      \bcd{(\lambd[x][N])}$,
      yaiku $\lambd[x][N'] \bredpar
      \lambd[x][\bcd{N}]$
      miturut \olref{defn:bcd2}.
      Iki ndherek saka
      \olref{defn:bredpar2} lan
      hipotesis induksi.
    \item Yen aturan pungkasan
      \olref{defn:bredpar3}:
      $M$ awujud $PQ$ lan
      $M'$ awujud $P'Q'$
      kanggo sawetara $P$,
      $Q$, $P'$, lan $Q'$,
      kanthi $P \bredpar P'$
      lan $Q \bredpar Q'$.
      Miturut hipotesis induksi,
      $P' \bredpar \bcd{P}$
      lan $Q' \bredpar \bcd{Q}$.
      \begin{enumerate}
        \item Yen $P$ awujud
          $\lambd[x][N]$ kanggo
          sawetara $x$ lan $N$,
          $P'$ mesthi awujud
          $\lambd[x][N']$
          kanggo sawetara $N'$
          kanthi $N \bredpar N'$.
          Miturut hipotesis induksi,
          $N' \bredpar \bcd{N}$
          lan $Q' \bredpar \bcd{Q}$.
          Mula $(\lambd[x][N'])Q' \bredpar
          \Subst{\bcd{N}}{\bcd{Q}}{x}$
          miturut \olref{defn:bredpar4}.
        \item Yen $P$ dudu
          abstraksi-$\lambd$,
          mula $P'Q' \bredpar
          \bcd{P}\bcd{Q}$
          miturut \olref{defn:bredpar3},
          lan sisih tengen iku
          $\bcd{PQ}$ miturut
          \olref{defn:bcd3}.
      \end{enumerate}
    \item Yen aturan pungkasan
      \olref{defn:bredpar4}: $M$
      awujud $(\lambd[x][N])Q$
      lan $M'$ awujud
      $\Subst{N'}{Q'}{x}$
      kanggo sawetara $x$,
      $N$, $Q$, $N'$, lan~$Q'$,
      kanthi $N \bredpar N'$
      lan $Q \bredpar Q'$.
      Miturut hipotesis induksi,
      $N' \bredpar \bcd{N}$
      lan $Q' \bredpar \bcd{Q}$.
      Miturut \olref{lem:comp},
      $\Subst{N'}{Q'}{x} \bredpar
      \Subst{\bcd{N}}{\bcd{Q}}{x}$;
      sisih tengen iku
      $\bcd{((\lambd[x][N])Q)}$.
  \end{enumerate}
\end{proof}

\begin{prob}
  Rampungna bukti \olref[lam][cr][pb]{lem:cont}.
\end{prob}

\begin{thm}\ollabel{thm:cr}
  $\bredpar$ nduweni
  sipat Church--Rosser.
\end{thm}

\begin{proof}
  Langsung saka \olref{lem:cont}.
\end{proof}

\end{document}
